\item Un estanque de agua a $0^\circ\text{C}$ se cubre con una capa de hielo de 4.00 cm de grueso. Si la temperatura del aire permanece constante a $-10.0^\circ\text{C}$, ¿qué intervalo de tiempo se requiere para que el espesor del hielo aumente a 8.00 cm? 
Sugerencia: Use la ecuación 7.18 en la forma
$$ \frac{dQ}{dt} = kA \frac{\Delta T}{x} $$
y advierta que el incremento de energía $dQ$ extraída del agua a través del espesor $x$ de hielo es la cantidad requerida para congelar un espesor $dx$ de hielo. Es decir, $dQ = L \rho A dx$, donde $\rho$ es la densidad del hielo, $A$ es el área y $L$ es el calor latente de fusión.
